\section{Core Modules}
\label{sec:manual-core-modules}

This section describes the core modules that are part of the CubedOS distribution. These modules
are intended to be used by most CubedOS applications. They provide basic services. Additional
information about a core module can be found in the comments in the module's top level package
and |API| package specifications, or in the MXDR file used by Mercury to generate the modules
|API| package.

\subsection{Log Server}
\label{sec:manual-log-server}

The log server is a core module that provides a simple logging facility. The default log server
outputs log messages to the console via Ada's standard library package |Ada.Text_IO|. For
systems where that package is not available, developers can write their own body for package
|CubedOS.Log_Server.Messages| that handles incoming messages in whatever way makes sense for
their application.

Log messages are intended to be short strings of human-readable text. Each message has an
associated log level, described below, and is automatically timestamped by the log server with a
boot-relative time in seconds.

The log server's API package provides a convenience procedure named |Log_Message| that takes
care of the details of constructing a message and sending it to the log server. It is expected
that most uses of the Log Server will be by way of calling |Log_Message|. Currently, the Log
Server does not send any reply messages. Logging is assumed to always be successful. This is an
area for future improvement.

Of particular interest is the |Log_Level| parameter to |Log_Message|. This parameter is used to
indicate the severity of the log message. It is an enumeration type. The Log levels are intended
to follow those used by the Unix syslog facility, both in name and meaning. The following log
levels are defined:

\begin{itemize}
\item |Debug|: Messages that are only of interest to developers. These messages are typically
very detailed and are used to diagnose problems in the system. Debug messages are not typically
enabled in a production system.

\item |Informational|: Messages that are of interest to users. These messages are typically
high-level and are used to inform the user of the system's normal operation.

\item |Notice|: Like informational messages, notice messages occur during normal system
operation. However, they log particularly significant events that are meant to stand out
compared to ordinary informational messages.

\item |Warning|: Messages that indicate a potential problem. The system is still functioning
correctly, but the message indicates that something might be wrong.

\item |Error|: Messages that indicate that a problem has occurred. The system is still
functioning, but the message indicates that something unintended or unexpected has happened.

\item |Critical|: Critical messages are similar to error messages, except with more severity and
criticality. A critical message indicates that the system is in a bad state and that the problem
should be addressed soon if the system is to continue functioning normally.

\item |Alert|: Messages that indicate a problem that has occurred that is so severe that the
system cannot function normally without corrective action. The system will likely soon fail
completely.

\item |Emergency|: Emergency messages are the most severe. They indicate that the system has
failed, is unusable, or soon will be, and that immediate action is required.
\end{itemize}

The timestamps added to every log message show the time the message was processed by the log
server relative to when the system started. Timestamps are in seconds. This time includes the
overhead of passing the message from the sending module to the log server. Because multiple
messages might be queued in the log server's mailbox, and because of task priority issues, the
timestamp might indicate a time that is, potentially, significantly later than when the message
was generated by the sending module.

One area for future work in the log server would be for the sending module to generate a
timestamp immediately before when the message is sent and pass that timestamp to the log server
as part of the message. In that case, the timestamps would be a more accurate reflection of when
the message was generated, being free from queuing and task priority effects. The convenience
procedure |Log_Message| could take care of the timestamp management details, so the developer
wouldn't need to do any extra work.

Of course, many applications won't care about such precise timestamps. The current
implementation will likely be sufficient for those applications. However, applications with very
tight timing requirements might benefit from better debugging and monitoring that would be made
possible by more accurate timestamps.

An implementation that displayed timestamps using an absolute time was considered. However, such
an implementation would need access to a real-time clock via, for example, package
|Ada.Calendar|. While |Ada.Calendar| is supported under the Jorvik profile, some small embedded
environments nevertheless do not support it. To keep CubedOS as widely applicable as possible,
it was decided to avoid the use of |Ada.Calendar| and instead focus on only boot-relative
timestamps.

Other future enhancements to the log server include (but are not limited to) the following:
\begin{itemize}
\item There should be a way to specify the log level that is displayed. Under normal operation,
users will not want to see Debug messages, for example. Furthermore, because such messages might
be plentiful and could potentially consume significant resources to transmit in volume, the
filtering of those messages should be done in |Log_Message| and not by the log server itself.
That would avoid transmitting messages that are of no interest.

\item As an extension of the idea above, it should be possible to control the log level based on
the module ID of the sending module. The use-case for this ability is when the developer would
like to log Debug messages from a particular module being debugged without having to see all
Debug messages generated from all other modules too.
\end{itemize}
